\documentclass[handout]{beamer}
\mode<presentation>
\usetheme[secheader]{Boadilla}
\usecolortheme{whale}

\usepackage[utf8]{inputenc}
\usepackage[T1]{fontenc}
\usepackage[indonesian]{babel}
\usepackage{lmodern}
\usepackage{graphicx}
\usepackage{bookmark}

\setbeamertemplate{footline}
{
	\leavevmode%
	\hbox{%
		\begin{beamercolorbox}[wd=.333333\paperwidth,ht=2.25ex,dp=1ex,center]{author in head/foot}%
			\usebeamerfont{author in head/foot}\insertshortauthor%
		\end{beamercolorbox}%
		\begin{beamercolorbox}[wd=.433333\paperwidth,ht=2.25ex,dp=1ex,center]{title in head/foot}%
			\usebeamerfont{title in head/foot}\insertshorttitle
		\end{beamercolorbox}%
		\begin{beamercolorbox}[wd=.233333\paperwidth,ht=2.25ex,dp=1ex,right]{date in head/foot}%
			\usebeamerfont{date in head/foot}\insertshortdate{}\hspace*{2em} \insertframenumber{} / \inserttotalframenumber\hspace*{2ex}
	\end{beamercolorbox}}%
	\vskip0pt%
}

\title[Bab 11 -- Library/Framework]{Pemrograman Web}
\subtitle{Bab 11: Library/Framework Front-End (jQuery/Bootstrap/Vue/React) untuk UI Dinamis}
\author{Cecep Suwanda, S.Si., M.Kom.}
\institute{Universitas Bale Bandung}
\date{\today}

\begin{document}

% Slide Judul
\begin{frame}
	\titlepage
\end{frame}

% Outline dan CPMK
\begin{frame}{Outline Bab 11}
	\begin{block}{Sub-CPMK 2.4}
		Menggunakan library atau framework front-end untuk antarmuka dinamis.
	\end{block}
	\begin{itemize}
		\item Timer: setTimeout, setInterval
		\item Asynchronous: Callback, Promise, async/await
		\item ES6+: Destructuring, Spread, Modul
		\item Class dan Pewarisan
		\item RegExp, Fetch API
		\item Closure, Prototype, Pola Desain
		\item Debugging, Unit Test, Performa
		\item Keamanan: XSS dan Sanitasi
		\item Library \& Framework Front-End
		\item Tooling: Bundler
	\end{itemize}
\end{frame}

% ========== SECTION 11-1: Timer ==========
\begin{frame}{Timer: setTimeout dan setInterval}
	\begin{itemize}
		\item \texttt{setTimeout(callback, ms)}: jalankan callback \textbf{sekali} setelah jeda.
		\item \texttt{setInterval(callback, ms)}: jalankan callback \textbf{berulang} setiap interval.
		\item Pembatalan: \texttt{clearTimeout(id)}, \texttt{clearInterval(id)}.
		\item \texttt{requestAnimationFrame(fn)}: sebelum repaint (60fps).
		\item Timer bersifat \textbf{asinkron}---kode setelah \texttt{setTimeout} dieksekusi tanpa menunggu.
	\end{itemize}
	\begin{block}{Debounce \& Throttle}
		\textbf{Debounce}: tunggu event berhenti, jalankan sekali (search-as-you-type). \textbf{Throttle}: maksimal sekali per interval (scroll handler).
	\end{block}
\end{frame}

% ========== SECTION 11-2: Asynchronous ==========
\begin{frame}{Pengenalan Asynchronous: Callback}
	\begin{itemize}
		\item Kode \textbf{asinkron} tidak menunggu operasi selesai; hasil ditangani lewat \textbf{callback}.
		\item JavaScript \textit{single-threaded}: satu operasi per waktu.
		\item Masalah: \textbf{callback hell} (pyramid of doom)---callback bersarang sulit dibaca.
		\item Solusi: Promise dan async/await (sintaks lebih linier).
	\end{itemize}
	\begin{block}{Event Loop}
		Call stack + task queue. \texttt{setTimeout(fn, 0)} masuk antrian, menunggu call stack kosong.
	\end{block}
\end{frame}

% ========== SECTION 11-3: ES6+ ==========
\begin{frame}{ES6+: Destructuring, Spread, Modul}
	\begin{itemize}
		\item \textbf{Destructuring}: \texttt{const [a,b]=[1,2]; const \{name\}=user;}
		\item \textbf{Spread}: \texttt{[...arr]}, \texttt{\{...obj\}}; \textbf{Rest}: \texttt{...args} di parameter.
		\item \textbf{ES6 Modules}: \texttt{import}/\texttt{export}; scope sendiri, strict mode; \texttt{<script type="module">}.
		\item Named export vs default export; gunakan named untuk utility, default untuk komponen utama.
	\end{itemize}
\end{frame}

% ========== SECTION 11-4: Class ==========
\begin{frame}{Class dan Pewarisan (ES6)}
	\begin{itemize}
		\item \texttt{class} mendefinisikan constructor dan method; \texttt{constructor()} dipanggil saat \texttt{new}.
		\item Pewarisan: \texttt{extends} dan \texttt{super()} untuk memanggil parent.
		\item \texttt{class} = syntactic sugar di atas prototype-based inheritance.
		\item \texttt{static} membuat method milik class (bukan instance).
	\end{itemize}
\end{frame}

% ========== SECTION 11-5: RegExp ==========
\begin{frame}{Regular Expression (RegExp) Dasar}
	\begin{itemize}
		\item Regex: pola teks untuk pencarian, validasi, penggantian.
		\item Metode: \texttt{regex.test(str)}, \texttt{str.match(regex)}, \texttt{str.replace(regex, repl)}.
		\item Pola: \texttt{.} (karakter apapun), \texttt{\textbackslash d} (digit), \texttt{\textbackslash w} (word), \texttt{+}/\texttt{*}, \texttt{\^{}}/\texttt{\$}, \texttt{[abc]}.
		\item Validasi format (email, telepon, password); server-side check tetap wajib.
	\end{itemize}
\end{frame}

% ========== SECTION 11-6: Promise ==========
\begin{frame}{Promise dan async/await}
	\begin{itemize}
		\item \textbf{Promise}: nilai di masa depan; state: pending, fulfilled, rejected.
		\item \texttt{.then()}, \texttt{.catch()}, \texttt{.finally()}; chain-able.
		\item \textbf{async/await}: sintaks linier; \texttt{await} menunggu Promise; error dengan \texttt{try/catch}.
		\item async/await direkomendasikan---lebih mudah dibaca daripada chain \texttt{.then()}.
	\end{itemize}
\end{frame}

% ========== SECTION 11-7: Fetch ==========
\begin{frame}{Fetch API}
	\begin{itemize}
		\item \texttt{fetch(url)} mengembalikan Promise; \texttt{res.json()} juga Promise.
		\item Mendukung GET, POST, PUT, DELETE via objek opsi.
		\item \textbf{Penting}: \texttt{fetch()} hanya reject pada network error; 404/500 tetap resolve---periksa \texttt{res.ok} atau \texttt{res.status}.
		\item Alur: loading indicator $\to$ fetch $\to$ parse JSON $\to$ render DOM $\to$ error handling.
	\end{itemize}
\end{frame}

% ========== SECTION 11-8: Closure ==========
\begin{frame}{Closure dan Lexical Scope}
	\begin{itemize}
		\item \textbf{Closure}: fungsi ``mengingat'' variabel dari scope luar setelah fungsi luar selesai.
		\item \textbf{Lexical scoping}: scope ditentukan posisi kode saat ditulis.
		\item Pola: factory function, data privacy (counter privat), event handler.
		\item Hati-hati loop + \texttt{var} + closure; gunakan \texttt{let} (block-scoped).
	\end{itemize}
\end{frame}

% ========== SECTION 11-9: Prototype ==========
\begin{frame}{Prototype dan Pewarisan}
	\begin{itemize}
		\item Setiap objek punya \textbf{prototype}; pewarisan lewat \textbf{prototype chain}.
		\item \texttt{\_\_proto\_\_} / \texttt{Object.getPrototypeOf()}; \texttt{F.prototype} untuk constructor.
		\item Pencarian: objek $\to$ prototype $\to$ ... $\to$ \texttt{Object.prototype} $\to$ null.
		\item \texttt{class extends} adalah sugar di atas prototype chain.
	\end{itemize}
\end{frame}

% ========== SECTION 11-10: Pola Desain ==========
\begin{frame}{Pola Desain: Module \& Singleton}
	\begin{itemize}
		\item \textbf{Module pattern}: IIFE atau closure untuk scope privat; hanya API publik di-return.
		\item ES6 modules sudah native; IIFE masih relevan untuk konsep.
		\item \textbf{Singleton}: hanya satu instance (konfigurasi, state global).
		\item Pola desain = solusi terbukti, bukan aturan wajib; gunakan jika masalah sesuai.
	\end{itemize}
\end{frame}

% ========== SECTION 11-11: Debugging ==========
\begin{frame}{Debugging dan DevTools}
	\begin{itemize}
		\item Panel: \textbf{Console} (log, error), \textbf{Sources} (breakpoint, step-through), \textbf{Network}, \textbf{Elements}.
		\item Teknik: \texttt{console.log}, \texttt{console.table}, \texttt{debugger;}, breakpoint, watch expression.
		\item \textbf{Breakpoint > console.log}: lihat call stack, variabel, step-by-step.
		\item Conditional breakpoint untuk loop besar.
	\end{itemize}
\end{frame}

% ========== SECTION 11-12: Unit Test ==========
\begin{frame}{Pengujian Dasar: Unit Test}
	\begin{itemize}
		\item Unit test: periksa fungsi/modul secara terisolasi; otomatis, repeatable.
		\item Framework: Jest, Vitest, Mocha.
		\item Pola AAA: Arrange-Act-Assert.
		\item Fokus: happy path, edge cases, error cases. 100\% coverage bukan jaminan.
	\end{itemize}
\end{frame}

% ========== SECTION 11-13: Performa ==========
\begin{frame}{Performa: Praktik Baik}
	\begin{itemize}
		\item Bottleneck: DOM manipulation berlebihan (layout thrashing), event tanpa debounce/throttle.
		\item Praktik: batch DOM (\texttt{DocumentFragment}), debounce/throttle, \texttt{requestAnimationFrame}, lazy load, profiling.
		\item \textbf{Layout thrashing}: baca layout, tulis DOM bergantian $\to$ forced sync layout. Solusi: kumpulkan baca dulu, lalu tulis.
		\item Web Workers untuk komputasi berat.
	\end{itemize}
\end{frame}

% ========== SECTION 11-14: Keamanan ==========
\begin{frame}{Keamanan: XSS dan Sanitasi}
	\begin{itemize}
		\item \textbf{XSS}: input berisi script disisipkan tanpa sanitasi $\to$ dieksekusi di browser korban.
		\item Pencegahan: gunakan \texttt{textContent} (bukan \texttt{innerHTML}); escape output; CSP header; validasi server-side.
		\item \textbf{Prinsip}: jangan percaya input pengguna. Validasi dan sanitasi di client dan server.
	\end{itemize}
\end{frame}

% ========== SECTION 11-15: Library/Framework ==========
\begin{frame}{Library dan Framework Front-End}
	\begin{itemize}
		\item \textbf{Library}: Anda memanggil; \textbf{Framework}: framework memanggil Anda (inversion of control).
		\item React (komponen, virtual DOM), Vue (mudah, reaktif), Angular (TypeScript, enterprise).
		\item jQuery (legacy; hindari untuk proyek baru).
		\item Component-based architecture: UI dipecah komponen kecil.
		\item \textbf{Kuasai JavaScript dulu} sebelum framework.
	\end{itemize}
\end{frame}

% ========== SECTION 11-16: Tooling ==========
\begin{frame}{Alat Modern: Bundler \& Tooling}
	\begin{itemize}
		\item \textbf{Bundler}: gabung file JS/CSS/gambar $\to$ optimal untuk produksi. Vite, webpack, Parcel.
		\item \textbf{Vite}: \texttt{npm create vite@latest}; cepat, HMR, ES modules; pilihan standar proyek baru.
		\item npm/pnpm, ESLint, Prettier, TypeScript, Git.
	\end{itemize}
\end{frame}

% ========== SECTION 11-17: Praktik Proyek ==========
\begin{frame}{Praktik Proyek}
	\begin{itemize}
		\item Ide: Todo List, Weather App, Quiz App, Form Registrasi, Dashboard.
		\item Metodologi: user stories $\to$ wireframe $\to$ struktur data $\to$ implementasi bertahap $\to$ testing $\to$ refactoring.
		\item Konsep: DOM, event, Fetch, validasi, (opsional) library.
	\end{itemize}
\end{frame}

% ========== Rangkuman dan Penutup ==========
\begin{frame}{Rangkuman Bab 11}
	\begin{itemize}
		\item Timer, async (callback, Promise, async/await), Fetch API.
		\item ES6+ (destructuring, spread, modules), class, RegExp, closure, prototype.
		\item Debugging, unit test, performa, keamanan XSS.
		\item Library/framework: jQuery, Bootstrap, Vue, React.
		\item Kuasai JavaScript dasar dulu; framework adalah abstraksi.
	\end{itemize}
\end{frame}

\begin{frame}{}
	\centering
	\vfill
	\textbf{\Large Pertanyaan?}
	\vfill
\end{frame}

\end{document}
